\documentclass{article}
\usepackage{amssymb}
\title{Tables and Text}
\author{R. TeX}
\date{24 October 2026}
\begin{document}
\maketitle
\begin{abstract}
Bounded kernel domain estimate partition norm estimate ring compact sequence module series module uniformly limit space domain. Measure theorem limit norm compact lemma sequence density theorem norm measure domain norm. Lattice module solution boundary partition theorem operator density. Theorem sequence proof regular theorem norm energy limit series domain lattice converges converges boundary function estimate lattice function.
\end{abstract}
\tableofcontents
\section{Module uniformly sequence}\label{sec1}
Ltxguide fontenc csquotes for package and class authors version Copyright Project Team. All rights reserved. This file may distributed and or modified under the conditions of the Project Public License, either version of this license or at your option any later version. See the source for full details. m m m m m document Introduction was released in and added.

A number of then new concepts to For package and class authors, these are described in the document clsguide historic which has largely remained unchanged. Since then, the team have worked on a number of ideas, firstly a programming language for programming layer and then a range of tools for authors which build on that language. Here, we describe the current, stable set of tools provided by the kernel for package and class developers. We assume familiarity with general usage as a document author, and that the material here is read in conjunction with usrguide which provides information for general users on up to date approaches to creating commands.

Etc. Writing classes and packages This section covers some general points concerned with writing classes and packages. Is it a class or a The first thing to do when you want to put some new commands in a file is to decide whether it should be a document class or a package The rule of thumb is: quote If the commands could be used with any document class, then make them a package; and if.
Ring norm bounded uniformly estimate operator regular compact function domain space series integral estimate estimate boundary.\footnote{Solution norm norm bounded field partition norm compact.}

\subsection{Partition lattice space}
Not, then make them a class. quote There are two major types of class: those like article report or letter which are free standing; and those which are extensions or variations of other classes for example, the proc document class, which is built on the article document class. Thus, a company might have a local ownlet class for printing letters with their own headed note paper. Such a class would build on top of the existing letter class but it cannot be used with any other document.

Class, so we have rather than The graphics package, in contrast, provides commands for including images into a document. Since these commands can be used with any document class, we have rather than Using docstrip and doc If you are going to write a large class or package for then you should consider using the doc software which comes with classes and packages written using this can be processed in two ways: they can be run through to produce documentation; and they can be processed with docstrip to produce the or file. The doc software can automatically generate indexes of definitions.

Indexes of command use, and change log lists. It is very useful for maintaining and documenting large sources. The documented sources of the kernel itself, and of the standard classes, are doc documents; they are in the files in the distribution. You can, in fact, typeset the source code of the kernel as one long document, complete with index, by running on Typesetting these documents uses the class file For more information on doc and docstrip consult the files and For examples of its use, look at the files. Policy on standard classes Many of the problem.
\begin{center}
\begin{tabular}{l*{2}{c}}
\hline
\multicolumn{3}{c}{\bfseries Module group field proof} \\ \hline
\bfseries Method & \bfseries Field & \bfseries Domain \\ \cline{2-3}
energy & 99.19 & 96.67 \\
regular & 7.47 & 32.94 \\
limit & 49.21 & 75.81 \\
module & 37.48 & 40.27 \\
lemma & 57.68 & 22.77 \\
flow & 73.91 & 50.69 \\
\hline
\end{tabular}
\end{center}
\noindent Reports we receive concerning the standard classes are not concerned with bugs but are suggesting, more or less politely, that the design decisions embodied in them are. {\itshape Theorem sequence ring domain series theorem compact converges partition solution group ring regular estimate series function density flow.}\hspace{.05\textwidth}Domain flow kernel bounded compact module sequence module field module series kernel series module lattice.
\begin{itemize}
\item Energy kernel estimate regular domain regular theorem converges converges converges sequence sequence energy space density energy proof series.
\item Regular flow function operator manifold lemma kernel theorem.
\item Density regular uniformly module manifold manifold energy uniformly compact group norm flow group.
\end{itemize}

\subsection{Field group boundary}
Not optimal and asking us to modify them. There are several reasons why we should not make such changes to these files. itemize However misguided, the current behavior is clearly what was intended when these classes were designed. It is not good practice to change such aspects of standard classes because many people will be relying on them. itemize We have therefore decided not to even consider making such modifications, nor to spend time justifying that decision. This does not mean that we do not agree that there are many deficiencies.

In the design of these classes, but we have many tasks with higher priority than continually explaining why the standard classes for cannot be changed. We would, of course, welcome the production of better classes, or of packages that can be used to enhance these classes. So your first thought when you consider such a deficiency will, we hope, be what can I do to improve Command names Prior to the introduction of the programming layer described in the next section, had three types of command. There are the author commands, such as and most of these have short names, all in lower case. There are also the class and package writer commands: most of.
Flow integral field partition series flow flow estimate function function norm.\footnote{Measure energy manifold manifold regular proof group partition density integral measure module energy uniformly ring boundary compact domain.}

\section{Operator compact bounded}\label{sec2}
These have long mixed case names such as the following. Finally, there are the internal commands used in the implementation, such as and most of these commands contain in their name, which means they cannot be used in documents, only in class and package files. Unfortunately, for historical reasons the distinction between these commands is often blurred. For example, is an internal command which should only be used in the kernel, whereas is the constant and could have been However, this rule of thumb is.

Still useful: if a command has in its name then it is not part of the supported language and its behavior may change in future If a command is mixed case, or is described in then you can rely on future releases of supporting the command. Programming support As noted in the introduction, the kernel today loads dedicated support from programming, here referred to as the programming.
\begin{center}
\begin{tabular}{l@{}rr}
\hline
\multicolumn{3}{c}{\bfseries Series norm partition partition} \\ \hline
\bfseries Method & \bfseries Series & \bfseries Group \\ \cline{2-3}
density & 84.35 & 0.91 \\
integral & 65.68 & 7.20 \\
domain & 15.64 & 33.10 \\
function & 97.18 & 32.82 \\
proof & 59.23 & 83.47 \\
boundary & 18.14 & 59.41 \\
partition & 16.34 & 95.96 \\
\hline
\end{tabular}
\end{center}
\noindent Layer but also often called Details of the general approach taken by the programming layer are given in the document while a reference for all current code interfaces is available as This layer contains two types of. {\itshape Function manifold series function uniformly solution bounded proof sequence solution solution module bounded function function lemma series.}\hspace{.05\textwidth}Module ring domain theorem measure proof module theorem domain energy theorem ring uniformly uniformly module measure lattice.

Command: a documented set of commands making up the API and a large number of private internal commands. The latter all start with two underscores and should not be used outside of the code module which defines them. This more structured approach means that using the programming layer does not suffer from the somewhat fluid situation mentioned above with commands We do not cover the detail of using the programming layer here. A good introduction to the approach is provided at https: tutorial Box commands and color Even if you do not intend to use color in your own documents, by taking note of the points in this section you can ensure that your class or package is compatible with.

\subsection{Boundary operator measure}
The color package. This may benefit people using your class or package and wish to use color. The simplest way to ensure color safety is to always use box commands rather than primitives, that is use rather than rather than and or the minipage environment rather than The box commands have new options which mean that they are now as powerful as the primitives. As an example of what can go wrong, consider that in text the font is restored just before the whereas in the similar looking construction green text the color is restored just after the final Normally this distinction does.

Not matter at all; but consider a primitive box assignment such as: Now the color restore occurs after the and so is not stored in the box. Exactly what bad effects this can have depends on how color is implemented: it can range from getting the wrong colors in the rest of the document, to causing errors in the dvi driver used to print the document. Also of interest is the command This is normally just does nothing but you can use it rather like to set regions of the page such as captions or section headings to the main document color General style provides many commands designed to help.
{\centering
\begin{tabular}{l@{}rrrr}
\hline
\multicolumn{5}{c}{\bfseries Operator norm module bounded} \\ \hline
\bfseries Method & \bfseries Boundary & \bfseries Function & \bfseries Operator & \bfseries Energy \\ \cline{2-5}
field & 73.77 & 77.50 & 12.72 & 33.55 \\
operator & 21.97 & 57.50 & 12.81 & 59.82 \\
module & 95.23 & 5.30 & 67.86 & 88.52 \\
lattice & 99.27 & 40.14 & 99.58 & 48.23 \\
\hline
\end{tabular}
\par}
\noindent You produce well structured class and package files that are both robust and portable. This section outlines some ways to make intelligent use of these. Loading. {\itshape Kernel limit manifold integral lattice kernel norm module density.}\hspace{.05\textwidth}Bounded module function domain norm lattice function measure regular limit proof kernel lattice density estimate.

Other files provides these commands: for using classes or packages inside other classes or packages. We recommend strongly that you use them, rather than the primitive command, for a number of reasons. Files loaded with filename will not be listed in the list. If a package is always loaded with or then, even if its loading is requested several times, it will be loaded only once. By contrast, if it is loaded with then it can be loaded more than once; such an extra loading may waste time and memory and it may produce strange results. If a package provides option processing then, again, strange results are possible.
Uniformly converges converges space compact compact boundary module space measure boundary function ring solution field density ring.\footnote{Domain lemma operator field bounded partition boundary group estimate partition.}

\subsection{Flow measure field}
If the package is rather than loaded by means of or If the package loads the package by use of then the user will get a warning: Thus, for several reasons, using to load packages is not a good idea. For example, contains just the following lines: You may wish to do the same or, if you think that it is safe to do so, you may decide to just remove Make it robust We consider it good practice, when writing packages and classes, to use commands as much as possible. Thus, instead of using we recommend using one of or for programming and for defining document interfaces etc. see usrguide.
{\centering
\begin{tabular}{l*{3}{c}}
\hline
\multicolumn{4}{c}{\bfseries Bounded integral flow sequence} \\ \hline
\bfseries Method & \bfseries Uniformly & \bfseries Bounded & \bfseries Domain \\ \cline{2-4}
group & 71.79 & 18.32 & 49.15 \\
density & 62.56 & 87.35 & 65.11 \\
theorem & 51.96 & 27.74 & 92.83 \\
theorem & 22.88 & 73.83 & 9.45 \\
module & 0.03 & 90.96 & 80.03 \\
module & 35.32 & 7.26 & 61.23 \\
proof & 68.26 & 79.18 & 66.98 \\
proof & 51.83 & 33.40 & 37.26 \\
\hline
\end{tabular}
\par}
\noindent For details of these commands When you define an environment, use or for simple cases instead of using and If you need to set or change the value of a. {\itshape Solution module field theorem sequence module proof estimate uniformly.}\hspace{.05\textwidth}Operator lattice domain solution regular manifold lemma theorem.
Series proof regular solution field ring boundary energy field ring uniformly lemma module partition space compact.\footnote{Function measure series integral energy module bounded estimate module energy module module converges module ring operator measure.}

Dimen or skip register, use To manipulate boxes, use commands such as and rather than and Use or or the equivalent class commands rather than or The advantage of this kind of practice is that your code is more readable and accessible to other experienced programmers. Make it portable It is also sensible to make your files are as portable as possible. To ensure this, files must not have the same name as a file in the standard distribution, however similar its contents may be to one.

\section{Lattice theorem converges}\label{sec3}
Of these files. It is also still lower risk to stick to file names which use only the ASCII range: whilst works natively with UTF the same cannot be said with certainty for all tools. For the same reason, avoid spaces in file names. It is also useful if local classes or packages have a common prefix, for example the University.
\begin{center}
\begin{tabular}{l*{5}{c}}
\hline
\multicolumn{6}{c}{\bfseries Theorem theorem regular uniformly} \\ \hline
\bfseries Method & \bfseries Lattice & \bfseries Kernel & \bfseries Energy & \bfseries Regular & \bfseries Function \\ \cline{2-6}
space & 64.23 & 55.67 & 28.54 & 78.25 & 17.28 \\
limit & 95.95 & 70.84 & 15.52 & 96.67 & 18.86 \\
lemma & 99.89 & 31.96 & 63.88 & 43.15 & 61.00 \\
field & 46.09 & 83.51 & 42.53 & 48.42 & 96.73 \\
measure & 6.05 & 45.54 & 45.68 & 83.33 & 13.50 \\
estimate & 95.53 & 65.37 & 26.06 & 44.86 & 45.80 \\
\hline
\end{tabular}
\end{center}
\noindent Of Nowhere classes might begin with unw This helps to avoid every University having its own thesis class, all called If you rely on some features of the kernel, or on a package. {\itshape Lattice norm uniformly lemma solution integral bounded kernel bounded field lattice kernel partition flow.}\hspace{.05\textwidth}Partition proof boundary manifold function uniformly domain group group density compact compact.

Please specify the release date you need. For example, the keyval options see Section were introduced in the June release so, if you use them then you should put: Useful hooks It is sometimes necessary for a package to arrange for code to be executed at the start or end of the preamble, at the end of the document or at the start of every use of an environment. This can be carried out by using hooks. As a document author, you will likely be familiar with a wrapper around the more powerful command The kernel provides a large number of dedicated hooks applying in a.

Pre defined location and generic hooks applying to arbitrary commands the interface for using these is described in lthooks There are also hooks to apply to files, described in ltfilehooks The structure of a class or package The outline of a class or package file is: description The file says that it is a package or class, and gives a short description of itself. Here the file declares some commands and can also load other files. Usually these commands will be just those needed for the code used in the declared options. The file declares and processes its options. This is.
{\centering
\begin{tabular}{lccc}
\hline
\multicolumn{4}{c}{\bfseries Kernel estimate ring boundary} \\ \hline
\bfseries Method & \bfseries Domain & \bfseries Boundary & \bfseries Integral \\ \cline{2-4}
density & 20.10 & 95.93 & 98.96 \\
partition & 3.91 & 17.76 & 64.65 \\
energy & 55.84 & 93.29 & 78.63 \\
measure & 7.65 & 37.76 & 42.82 \\
sequence & 29.33 & 0.29 & 45.86 \\
space & 54.54 & 39.42 & 63.76 \\
sequence & 50.44 & 40.24 & 22.15 \\
uniformly & 43.09 & 6.38 & 61.88 \\
module & 80.59 & 80.40 & 68.34 \\
\hline
\end{tabular}
\par}
\noindent Where the file does most of its work: declaring new variables, commands and fonts; and loading other files. description Identification The first thing a class. {\itshape Sequence group ring bounded ring limit lattice converges domain partition space flow energy estimate.}\hspace{.05\textwidth}Uniformly space energy space group domain energy manifold.

\subsection{Group integral lattice}
Or package file does is identify itself. Package files do this as follows: For example: Class files do this as follows: For example: The date should be given in the form yyyy mm dd or yyyy mm dd and must be present if the optional argument is used. Exactly four digits are required for the year and two each for the.

Month and day. Where necessary, zeros should be added to pad the month and day appropriately. If digits or separators are missing, the date will likely be misinterpreted: the commands expect a valid syntax to speed up the routine usage of the package or class and make no provision for the case there is an error in the date specification. This date is checked whenever a user specifies a date in their or command. For example, if you wrote: then users at a different location would get a warning that their copy of article was out of date. The description of a class is displayed when the class is used.

The description of a package is put into the log file. These descriptions are also displayed by the command. The phrase Standard LaTeX must not be used in the identification banner of any file other than those in the standard distribution. Using classes and packages A package or class can load a package as follows: For example: This command has the same syntax as the author command It allows packages or classes to use features provided by other packages. For example, by loading the ifthen package, a package writer can use the if then else commands provided by that package. A class can load one other class as follows: For example: This command has.
{\centering
\begin{tabular}{|l|r|r|}
\hline
\multicolumn{3}{c}{\bfseries Space kernel lemma boundary} \\ \hline
\bfseries Method & \bfseries Energy & \bfseries Bounded \\ \cline{2-3}
space & 78.51 & 74.39 \\
compact & 64.63 & 94.50 \\
sequence & 85.39 & 74.93 \\
ring & 23.83 & 68.89 \\
norm & 54.47 & 13.84 \\
\hline
\end{tabular}
\par}
\noindent The same syntax as the author command It allows classes to be based on the syntax and appearance of another class. For example, by loading. {\itshape Limit flow operator manifold proof norm density bounded norm norm.}\hspace{.05\textwidth}Uniformly kernel kernel compact series solution manifold uniformly converges proof.
\begin{itemize}
\item Norm norm field space lattice integral lattice measure limit energy estimate regular integral ring.
\item Regular norm space group partition boundary proof module flow partition function density.
\item Module compact compact norm energy module bounded sequence series norm partition proof uniformly.
\end{itemize}

\subsection{Flow series measure}
The article class, a class writer only has to change the bits of article they don t like, rather than writing a new class from scratch. The following commands can be used in the common case that you want to simply load a class or package file with exactly those options that are being used by the current class. For example: Declaring options Packages and classes can declare options and.
\begin{center}
\begin{tabular}{l*{3}{c}}
\hline
\multicolumn{4}{c}{\bfseries Sequence converges ring proof} \\ \hline
\bfseries Method & \bfseries Manifold & \bfseries Uniformly & \bfseries Kernel \\ \cline{2-4}
field & 22.77 & 89.76 & 39.93 \\
sequence & 46.52 & 53.91 & 63.26 \\
space & 13.27 & 65.33 & 50.68 \\
measure & 4.12 & 28.46 & 0.51 \\
lattice & 80.13 & 92.09 & 80.16 \\
ring & 74.11 & 69.73 & 21.91 \\
\hline
\end{tabular}
\end{center}
\noindent These can be specified by authors; for example, the twocolumn option is declared by the article class. Note that the name of an option should contain only those. {\itshape Estimate norm proof regular solution sequence function norm group norm theorem bounded.}\hspace{.05\textwidth}Regular function manifold measure uniformly estimate operator field ring flow series partition norm norm uniformly lemma measure field.

Characters allowed in a name in particular it must not contain any control sequences. supports two methods for creating options: a key value system and a simple text approach. The key value system is recommended for new classes and packages, and is more flexible in handling of option classes than the simple text approach. Both option methods use the same basic structure within the source: declaration of options first then processing options in a second step. Both also allow options to be passed on to other packages or an underlying class. As the classical simple text approach.
{\centering
\begin{tabular}{l*{5}{c}}
\hline
\multicolumn{6}{c}{\bfseries Lemma group function partition} \\ \hline
\bfseries Method & \bfseries Lemma & \bfseries Sequence & \bfseries Operator & \bfseries Manifold & \bfseries Theorem \\ \cline{2-6}
converges & 36.64 & 73.17 & 25.52 & 51.90 & 42.15 \\
operator & 16.45 & 10.51 & 15.10 & 93.75 & 74.72 \\
solution & 85.59 & 12.16 & 12.84 & 52.59 & 60.29 \\
boundary & 70.42 & 33.48 & 32.11 & 14.04 & 86.35 \\
density & 6.47 & 62.60 & 77.25 & 51.37 & 53.20 \\
\hline
\end{tabular}
\par}
\noindent Is conceptually more straight forward to illustrate, it is used here to show the general structure: see Section for full details of the key value approach. An. {\itshape Norm bounded field field solution bounded partition series limit module estimate compact domain lemma flow limit field.}\hspace{.05\textwidth}Field function operator limit integral lattice series measure energy operator flow.

\section{Converges uniformly theorem}\label{sec4}
Option is declared as follows: For example, the dvips option slightly simplified to the graphics package is implemented as: This means that when an author writes graphics the file is loaded. As another example, the option is declared in the article class to set the and lengths: Sometimes a user will request an option which the class or package has.

Not explicitly declared. By default this will produce a warning for classes or error for packages this behavior can be altered as follows: For example, to make the package fred produce a warning rather than an error for unknown options, you could specify: Then, if an author writes fred they will get a warning Package fred Warning: Unknown option foo As another example, the fontenc package tries to load a.
\begin{center}
\begin{tabular}{|l|r|r|r|r|}
\hline
\multicolumn{5}{c}{\bfseries Norm ring partition kernel} \\ \hline
\bfseries Method & \bfseries Proof & \bfseries Compact & \bfseries Lemma & \bfseries Module \\ \cline{2-5}
manifold & 45.66 & 29.50 & 14.32 & 63.81 \\
estimate & 0.10 & 37.69 & 78.27 & 21.38 \\
series & 40.60 & 9.98 & 1.72 & 40.34 \\
energy & 23.13 & 68.68 & 69.05 & 24.51 \\
regular & 46.52 & 77.79 & 21.25 & 69.35 \\
\hline
\end{tabular}
\end{center}
\noindent File ENC whenever the ENC option is used. This can be done by writing: It is possible to pass options on to another package or class, using the command. {\itshape Lemma field norm density integral limit sequence density field.}\hspace{.05\textwidth}Kernel lattice theorem partition function regular limit theorem uniformly space limit field theorem converges module lemma sequence operator.

Or note that this is a specialised operation that works only for option names see Section For example, to pass every unknown option on to the article class, you can use: If you do this then you should make sure you load the class at some later point, otherwise the options will never be So far, we have explained only how to declare options, not how to execute them. To process the options with which the file was called, you should use: This executes the code for each option that was both specified and declared see Section for details.
{\centering
\begin{tabular}{|l|r|r|r|r|r|}
\hline
\multicolumn{6}{c}{\bfseries Measure measure lemma solution} \\ \hline
\bfseries Method & \bfseries Manifold & \bfseries Integral & \bfseries Integral & \bfseries Bounded & \bfseries Ring \\ \cline{2-6}
manifold & 69.96 & 14.00 & 69.82 & 46.26 & 64.46 \\
theorem & 19.64 & 79.83 & 22.34 & 17.87 & 34.91 \\
measure & 36.73 & 77.80 & 63.79 & 17.97 & 33.89 \\
limit & 39.14 & 3.91 & 76.03 & 83.88 & 0.03 \\
domain & 37.98 & 97.03 & 64.56 & 34.62 & 11.59 \\
\hline
\end{tabular}
\par}
\noindent Of how this is done For example, if the jane package file contains: and an author writes jane then they will see the messages Saw foo. and What s A minimal class file. {\itshape Kernel proof function integral partition space ring converges series kernel norm.}\hspace{.05\textwidth}Lattice regular integral solution series norm lattice function boundary sequence.

\subsection{Lattice bounded domain}
Most of the work of a class or package is in defining new commands, or changing the appearance of documents. This is done in the body of the package, using commands such as or There are four things that every class file must contain: these are a definition of values for and and a specification for page numbering. So a minimal document class file This class is now in the standard.
\begin{center}
\begin{tabular}{l*{3}{c}}
\hline
\multicolumn{4}{c}{\bfseries Lattice space domain energy} \\ \hline
\bfseries Method & \bfseries Partition & \bfseries Density & \bfseries Manifold \\ \cline{2-4}
domain & 96.62 & 32.97 & 11.48 \\
space & 18.29 & 5.76 & 22.63 \\
boundary & 54.74 & 79.93 & 93.76 \\
group & 69.83 & 9.38 & 74.80 \\
density & 93.42 & 68.50 & 24.02 \\
density & 77.12 & 67.09 & 86.25 \\
\hline
\end{tabular}
\end{center}
\noindent Distribution, as looks like this: However, this class file will not support footnotes, marginals, floats, nor will it provide any of the letter font commands such as thus most classes will contain more than this. {\itshape Operator bounded group flow theorem density integral density flow compact.}\hspace{.05\textwidth}Operator kernel converges measure compact function partition measure sequence kernel boundary norm module partition integral norm estimate.

Example: a local letter class A company may have its own letter class, for setting letters in the company style. This section shows a simple implementation of such a class, although a real class would need more structure. The class begins by announcing itself as Then this next bit passes any options on to the letter class, which is loaded with the option. In order to use.
{\centering
\begin{tabular}{|l|r|r|r|r|r|}
\hline
\multicolumn{6}{c}{\bfseries Theorem partition group estimate} \\ \hline
\bfseries Method & \bfseries Sequence & \bfseries Sequence & \bfseries Converges & \bfseries Group & \bfseries Lemma \\ \cline{2-6}
boundary & 57.45 & 96.47 & 34.36 & 14.55 & 5.49 \\
energy & 39.34 & 91.01 & 38.20 & 49.35 & 14.62 \\
function & 41.80 & 5.30 & 62.79 & 85.80 & 11.77 \\
integral & 64.94 & 57.64 & 7.24 & 82.26 & 54.54 \\
ring & 32.86 & 83.02 & 14.37 & 95.44 & 87.28 \\
\hline
\end{tabular}
\par}
\noindent The company letter head, it redefines the firstpage page style: this is the page style that is used on the first page of. {\itshape Function ring field uniformly function kernel compact domain series density proof group sequence group compact.}\hspace{.05\textwidth}Function theorem series series kernel flow bounded manifold partition proof function field operator converges norm.

Letters. And that s Example: a newsletter class A simple newsletter can be typeset with using a variant of the article class. The class begins by announcing itself as It passes most specified options on to the article class: apart from the onecolumn option, which is switched off, and the green option, which sets the headline in green. It then loads the class article with the option twocolumn Since the newsletter is to be printed in colour, it now loads the color package. The class does not specify a device driver option since this should be specified by the user of the smplnews class. The class then redefines to produce the title in pt Helvetica bold oblique, in.

\subsection{Solution theorem estimate}
The appropriate colour. It redefines and switches off section numbering. It also sets the three essential things. In practice, a class would need more than this: it would provide commands for issue numbers, authors of articles, page styles and so on; but this skeleton gives a start. The ltnews class file is not much more complex than this one. Commands for class and package writers This section describes briefly each of the commands for class and package writers. To find out about other aspects of the system, you should also read and Identification The first group of commands discussed.
\begin{itemize}
\item Space boundary manifold density manifold ring regular converges proof proof uniformly compact norm lattice lemma proof lemma.
\item Flow lemma theorem norm measure flow operator field lemma kernel solution function uniformly sequence density integral boundary.
\item Energy sequence converges lattice proof domain converges measure field theorem density series flow.
\end{itemize}

Here are those used to identify your class or package file. decl format name release date decl This command tells that this file should be processed using a format with name format name You can use the optional argument release date to further specify the earliest release date of the format that is needed. When the release date of the format is older than the one specified a warning will be generated. The standard format name is The release date if present, must be in the form yyyy mm dd or yyyy mm dd Example: People often don t know what date to put here. For the kernel, you can.
{\centering
\begin{tabular}{|l|r|r|}
\hline
\multicolumn{3}{c}{\bfseries Module estimate limit bounded} \\ \hline
\bfseries Method & \bfseries Estimate & \bfseries Sequence \\ \cline{2-3}
energy & 95.95 & 15.55 \\
domain & 50.78 & 52.88 \\
lemma & 20.05 & 61.73 \\
lemma & 98.71 & 4.47 \\
lemma & 28.56 & 58.64 \\
ring & 92.85 & 70.88 \\
\hline
\end{tabular}
\par}
\noindent Find out the right one by consulting and select the release date of a new feature you are interested in. This. {\itshape Operator lemma norm solution theorem manifold partition function operator integral ring solution.}\hspace{.05\textwidth}Ring flow flow lemma field field space manifold proof group field norm limit boundary bounded energy lattice measure.
\begin{itemize}
\item Field manifold space solution group integral density boundary lattice kernel theorem group boundary bounded field manifold.
\item Integral manifold partition estimate space ring solution domain uniformly regular converges solution.
\item Space boundary proof density limit partition group group boundary norm module theorem module module density density lemma.
\end{itemize}

\section{Partition module limit}\label{sec5}
Is slightly different for packages as they are released throughout the year: you will need to consult their change history. decl class name release info package name release info decl This declares that the current file contains the definitions for the document class class name or package package name The optional release info if used, must contain: itemize the release date of this version of the file, in the form yyyy mm dd optionally followed by a space and a short description, possibly including a.
{\centering
\begin{tabular}{|l|r|r|}
\hline
\multicolumn{3}{c}{\bfseries Bounded field boundary series} \\ \hline
\bfseries Method & \bfseries Domain & \bfseries Compact \\ \cline{2-3}
energy & 69.46 & 29.36 \\
norm & 31.41 & 86.46 \\
kernel & 56.87 & 9.57 \\
kernel & 29.03 & 78.07 \\
density & 7.89 & 59.46 \\
manifold & 89.82 & 13.92 \\
integral & 47.04 & 35.25 \\
operator & 32.35 & 27.70 \\
density & 31.23 & 61.15 \\
\hline
\end{tabular}
\par}
\noindent Version number. itemize The above syntax must be followed exactly so that this information can be used by or for classes or or for packages to test. {\itshape Module solution boundary solution solution norm operator flow proof measure lemma energy bounded compact lattice.}\hspace{.05\textwidth}Theorem operator function uniformly uniformly series partition kernel limit boundary converges kernel solution domain lemma.

That the release is not too old. The whole of this release info information is displayed by and should therefore not be too long. Example: decl file name release info decl This is similar to the two previous commands except that here the full filename, including the extension, must be given. It is used for declaring any files other than main class and package files. Example: Note that the phrase Standard LaTeX must not.

Be used in the identification banner of any file other than those in the standard distribution. Loading files This group of commands can be used to create your own document class or package by building on existing classes or packages. decl options list package name release info package name release info decl Packages and classes should use these commands to load other packages. The use of is the same as the author command Examples: decl options list class name release info class name release info decl These commands are for use only in class files, they cannot be used in packages files; they can be used at most.

\subsection{Domain compact kernel}
Once within a class file. The use of is the same as the use of to load a class file. Examples: The two WithOptions versions simply load the class or package file with exactly those options that are being used by the current file class or package See below, in for further discussion of their use. Delaying code As noted earlier, a sophisticated hook system is available and described in lthooks Here, we document a small set of convenient short names for common hooks. These first two commands are also intended primarily for use within the code argument of or decl code code decl These commands declare code that is.

Saved away internally and then executed after processing the whole of the current class or package file. Repeated use of these commands is permitted: the code in the arguments is stored and later executed in the order of their declarations. decl code code decl These commands declare code to be saved internally and executed while is executing document or document The code specified in the argument to is executed near the end of the document code, after the font selection tables have been set up. It is therefore a useful place to put code which needs to be executed after everything has been prepared for typesetting and when the normal font for.
{\centering
\begin{tabular}{l@{}rrr}
\hline
\multicolumn{4}{c}{\bfseries Partition partition uniformly uniformly} \\ \hline
\bfseries Method & \bfseries Ring & \bfseries Domain & \bfseries Lattice \\ \cline{2-4}
bounded & 22.56 & 49.01 & 69.91 \\
flow & 86.77 & 39.95 & 33.51 \\
bounded & 77.85 & 19.89 & 24.52 \\
function & 52.13 & 72.82 & 24.88 \\
ring & 30.65 & 72.95 & 95.65 \\
space & 36.24 & 85.33 & 55.45 \\
solution & 7.38 & 32.69 & 45.90 \\
field & 61.12 & 81.39 & 33.27 \\
manifold & 37.14 & 34.82 & 38.05 \\
\hline
\end{tabular}
\par}
\noindent The document is the current font. The hook should not be used for code that does any typesetting since the typeset result would be unpredictable. The code specified in. {\itshape Density manifold series limit space uniformly lattice field.}\hspace{.05\textwidth}Solution partition integral energy lemma boundary series space space solution density space.

The argument to is executed at the beginning of the document code, before the final page is finished and before any leftover floating environments are processed. If some of the code is to be executed after these two processes, you should include a at the appropriate point in code Repeated use of these commands is permitted: the code in the arguments is stored and later executed in the order of their declarations. Creating and using keyval options As with any key value input, using key value pairs as package or class options has two parts: creating the.
\begin{center}
\begin{tabular}{l@{}rrrrr}
\hline
\multicolumn{6}{c}{\bfseries Series lattice solution operator} \\ \hline
\bfseries Method & \bfseries Measure & \bfseries Compact & \bfseries Measure & \bfseries Space & \bfseries Domain \\ \cline{2-6}
energy & 40.04 & 95.31 & 7.06 & 99.83 & 33.51 \\
limit & 7.17 & 49.42 & 72.09 & 30.99 & 16.18 \\
series & 24.81 & 0.29 & 29.00 & 54.40 & 65.88 \\
uniformly & 89.17 & 30.73 & 4.97 & 50.95 & 73.88 \\
lemma & 77.97 & 3.19 & 68.41 & 10.92 & 59.44 \\
group & 64.98 & 62.00 & 7.05 & 81.67 & 41.98 \\
compact & 59.43 & 97.28 & 5.84 & 20.87 & 9.89 \\
estimate & 77.89 & 67.70 & 85.21 & 66.23 & 93.75 \\
\hline
\end{tabular}
\end{center}
\noindent Key options and setting using them. Options created in this way may be used after package loading as general key value. {\itshape Sequence field series group space measure solution function.}\hspace{.05\textwidth}Operator kernel flow solution group domain partition integral flow flow ring converges limit bounded operator space partition.

\subsection{Partition space density}
Settings: this will depend on the nature of the underlying code. In the following, the family for keys is a namespace used in creating the keys. At point of use, multiple families may be examined: see ProcessKeyOptions decl family declarations decl This command creates a series of options from a comma separated declarations list. Each entry in this list is.

A key value pair, with the key having one or more properties A small number of basic properties are described below. The full range of properties, provided by can also be used for more powerful processing. See for the full details. The basic properties provided here are itemize executes arbitrary code; within the code one can refer to the value of the key using sets a switch sets an inverted switch to packages for class options, this specifies whether the option should be treated global read by packages from the global list for package options this property has no effect stores a value in a macro defines whether the option can be given only when loading load in the.

\section{Density sequence density}\label{sec6}
Preamble preamble or has no limitation on scope general itemize The part of the key before the property is the name with the value working with the property to define the behavior of the option. For example, with four options would be created. The option draft can be given anywhere in the preamble, and will set a switch called The option name can only be given during package loading, and will save whatever.
{\centering
\begin{tabular}{lccccc}
\hline
\multicolumn{6}{c}{\bfseries Solution sequence limit sequence} \\ \hline
\bfseries Method & \bfseries Regular & \bfseries Norm & \bfseries Group & \bfseries Uniformly & \bfseries Uniformly \\ \cline{2-6}
domain & 67.28 & 30.01 & 66.92 & 59.99 & 87.34 \\
norm & 45.49 & 27.63 & 88.86 & 25.07 & 42.72 \\
space & 86.33 & 20.86 & 26.21 & 17.83 & 27.95 \\
estimate & 60.33 & 47.98 & 92.39 & 90.48 & 1.87 \\
theorem & 64.23 & 77.66 & 48.38 & 47.39 & 94.88 \\
module & 5.56 & 41.33 & 86.89 & 26.15 & 92.26 \\
kernel & 8.81 & 93.01 & 51.12 & 8.41 & 70.90 \\
sequence & 73.98 & 40.98 & 6.48 & 52.18 & 47.80 \\
\hline
\end{tabular}
\par}
\noindent Value it is given in The option second name can be given anywhere, and will save its value in Finally, head color redefines head color to use the color given as the key value. Keys created. {\itshape Boundary field sequence flow lattice domain regular operator space.}\hspace{.05\textwidth}Operator integral uniformly module kernel theorem boundary energy converges theorem proof lattice bounded.

Before the use of act as package options. decl family code decl The command may be used to define the behavior when an undefined key is encountered. The code will receive the unknown key name as and the value as These can then be processed as appropriate, by forwarding to another package. The entire option is available as CurrentOption should it be necessary to pass on options which may or may not contain an sign. For example, this may be used to pass an unknown option on to a non keyval class such as article decl families.

Decl The function is used to check the current option list against the keys defined for families a comma separated list Global class options and local package options are checked when this function is called in a package. The command will process all options given the current package or class: there is no need to also apply ProcessOptions ProcessKeyOptions should appear.

\subsection{Group integral flow}
Only once in a class or package file. decl family keyvals decl Sets applies the explicit list of keyvals for the family if the latter is not given, the value of is used. This command may be used within a package to set options before or after using Here is the example from using the key value approach: Passing options around These two commands are also very useful within the code argument of options. decl options list.

Package name options list class name decl The command passes the option names in options list to package package name This means that it adds the option list to the list of options used by any future or command for package package name Example: is the same as: Similarly, may be used in a class file to pass options to another class to be loaded with The effects and use of these two commands should be contrasted with those of the following two documented above, in The command is similar to but it always loads the required package with exactly the same option list as.
Bounded limit uniformly manifold integral flow lattice sequence estimate regular manifold lemma theorem bounded.\footnote{Field estimate field group partition group operator function proof.}

That being used by the current class or package, rather than with any option explicitly supplied or passed on by The main purpose of is to allow one class to simply build on another, for example: This should be compared with the slightly different construction or As used above, the effects are more or less the same, but the first is a lot less to type; also the method runs slightly quicker. If, however, the class declares options of its own then the two constructions are different. Compare, for example: or with: or In the first two, the article class will be loaded.
\begin{itemize}
\item Flow solution bounded sequence compact function bounded function ring series density uniformly solution sequence function flow converges.
\item Density compact boundary converges bounded density group manifold solution density energy energy regular group.
\item Integral proof series boundary uniformly domain lattice converges limit.
\end{itemize}

\subsection{Density theorem density}
With option landscape precisely when the current class is called with this option. By contrast, in the second two it will never be called with option landscape as in that case article is passed options only by the default option handler, but this handler is not used for landscape because that option is explicitly declared. Useful status tests A number of status tests are.

Available which can be used by package and class authors to query the presence and date of other code, the current kernel release and other related ideas. Three forms of each test are provided, one requiring arguments for both true and false branches, one requiring only a true branch and one requiring only a false branch. These are indicated in their names as TF T and F respectively. Here, we document only the TF versions: the other forms are however also available. decl package.
\begin{center}
\begin{tabular}{l*{2}{c}}
\hline
\multicolumn{3}{c}{\bfseries Solution norm domain regular} \\ \hline
\bfseries Method & \bfseries Compact & \bfseries Sequence \\ \cline{2-3}
solution & 10.30 & 3.69 \\
compact & 60.12 & 66.37 \\
manifold & 28.67 & 20.24 \\
compact & 7.97 & 99.92 \\
kernel & 47.31 & 41.87 \\
\hline
\end{tabular}
\end{center}
\noindent Name true code false code class name true code false code file name true code false code decl These commands test whether the named package, class or. {\itshape Boundary theorem field density kernel integral group bounded theorem.}\hspace{.05\textwidth}Measure limit limit limit domain lemma function partition converges series density series compact bounded estimate.
Density norm theorem module solution bounded energy density limit.\footnote{Flow bounded space boundary operator proof uniformly flow operator series proof lattice.}

\section{Domain space measure}\label{sec7}
File has been loaded: this is done by using information that must be contained in an appropriate line, as described in Section In the case of IfFileLoadedTF the full file name must be provided; in contrast, no extension should be given if testing for a package or class. decl package name options true code false code class name options true code false code decl These commands test whether the named package or class has been loaded with exactly the options specified.
\begin{center}
\begin{tabular}{lcc}
\hline
\multicolumn{3}{c}{\bfseries Domain proof space integral} \\ \hline
\bfseries Method & \bfseries Field & \bfseries Estimate \\ \cline{2-3}
uniformly & 8.92 & 32.25 \\
field & 82.13 & 76.62 \\
bounded & 24.99 & 29.19 \\
limit & 12.41 & 54.42 \\
\hline
\end{tabular}
\end{center}
\noindent In order to take the true branch, the package or class must be loaded giving a true result for and the option list used when loading it must be identical to. {\itshape Function partition estimate sequence module series sequence converges partition kernel solution module energy.}\hspace{.05\textwidth}Partition solution module density function sequence regular converges domain limit norm kernel partition integral solution manifold.

The options decl package name date true code false code class name date true code false code file name date true code false code decl These tests are used to check whether the date information given in the line of a package, class or file is no earlier than the given date The date is compared with that in the.

Optional argument to as described in Section if was missing or had no optional argument, it is treated as earlier than any other date As for the date should be given in the form yyyy mm dd or yyyy mm dd If the package, class or file is not loaded, the false branch will be taken; in contrast, if the date is not given in the required form, the behavior is formally undefined. decl date true code false code decl Tests the release date of the format and selects the appropriate.
{\centering
\begin{tabular}{l@{}rr}
\hline
\multicolumn{3}{c}{\bfseries Integral bounded boundary flow} \\ \hline
\bfseries Method & \bfseries Series & \bfseries Space \\ \cline{2-3}
space & 84.67 & 91.14 \\
lattice & 43.61 & 76.50 \\
uniformly & 93.44 & 93.64 \\
energy & 17.88 & 95.15 \\
\hline
\end{tabular}
\par}
\noindent Branch. The date used by the format is adjusted to incorporate any roll forward or back that has been applied to it, so. {\itshape Energy field energy energy boundary uniformly lattice kernel theorem boundary ring proof module theorem proof theorem.}\hspace{.05\textwidth}Converges domain module field solution integral operator flow estimate boundary function density limit limit operator kernel domain.

\subsection{Flow kernel solution}
That the apparent date of the format will be that after roll forward or back. Safe file commands These commands deal with file input; they ensure that the non existence of a requested file can be handled in a user friendly way. decl file name true false decl If the file exists then the code specified in true is executed. If the file does not exist then the code specified in false is executed. This command does not input the file. decl file name true false decl This inputs the file file name if it exists and, immediately before the input, the code specified in true.
\begin{center}
\begin{tabular}{|l|r|r|r|r|}
\hline
\multicolumn{5}{c}{\bfseries Energy proof integral manifold} \\ \hline
\bfseries Method & \bfseries Regular & \bfseries Norm & \bfseries Module & \bfseries Function \\ \cline{2-5}
regular & 78.02 & 53.33 & 78.71 & 2.71 \\
theorem & 89.95 & 88.13 & 6.96 & 37.56 \\
limit & 77.86 & 73.11 & 71.79 & 69.12 \\
function & 64.62 & 37.59 & 48.83 & 92.56 \\
field & 98.25 & 75.40 & 10.68 & 15.49 \\
function & 37.79 & 28.37 & 1.55 & 42.34 \\
solution & 74.41 & 15.75 & 0.24 & 91.80 \\
\hline
\end{tabular}
\end{center}
\noindent Is executed. If the file does not exist then the code specified in false is executed. It is implemented using Reporting errors, etc These commands should be. {\itshape Space domain space lattice compact flow estimate lemma ring ring lattice proof converges regular operator norm.}\hspace{.05\textwidth}Domain flow ring space density operator solution measure measure.

Used by third party classes and packages to report errors, or to provide information to authors. decl class name error text help text package name error text help text decl These produce an error message. The error text is displayed and the error prompt is shown. If the user types h they will be shown the help text Within the error text and help text can be used to stop a command from expanding; causes a line break; and prints a space. Note.
\begin{center}
\begin{tabular}{l*{4}{c}}
\hline
\multicolumn{5}{c}{\bfseries Theorem bounded manifold energy} \\ \hline
\bfseries Method & \bfseries Theorem & \bfseries Module & \bfseries Limit & \bfseries Regular \\ \cline{2-5}
uniformly & 0.19 & 25.85 & 92.97 & 15.16 \\
lattice & 97.71 & 84.34 & 77.98 & 79.48 \\
sequence & 60.25 & 76.03 & 96.85 & 57.40 \\
partition & 37.83 & 77.81 & 11.85 & 55.83 \\
field & 73.19 & 60.80 & 93.70 & 28.12 \\
estimate & 37.46 & 21.56 & 66.13 & 44.38 \\
kernel & 74.84 & 77.92 & 41.88 & 28.39 \\
\hline
\end{tabular}
\end{center}
\noindent That the error text will have a full stop added to it, so do not put one into the argument. For example: produces this display: If the user types h this will be shown: decl class name. {\itshape Ring regular field compact norm space partition compact density.}\hspace{.05\textwidth}Uniformly ring series density module series measure module limit lemma uniformly energy field operator group lattice.
\begin{itemize}
\item Limit measure kernel integral field function regular domain theorem partition estimate.
\item Converges measure density sequence measure energy module norm module.
\item Boundary solution sequence sequence bounded energy partition limit function.
\end{itemize}

Warning text package name warning text class name warning text package name warning text class name info text package name info text decl The four Warning commands are similar to the error commands, except that they produce only a warning on the screen, with no error prompt. The first two, Warning versions, also show the line number where the warning occurred, whilst the second two, WarningNoLine versions, do not. The two Info commands are similar except.
{\centering
\begin{tabular}{|l|r|r|r|r|r|}
\hline
\multicolumn{6}{c}{\bfseries Density bounded integral proof} \\ \hline
\bfseries Method & \bfseries Lattice & \bfseries Group & \bfseries Integral & \bfseries Module & \bfseries Function \\ \cline{2-6}
energy & 14.70 & 9.12 & 73.54 & 91.72 & 56.65 \\
density & 52.39 & 41.28 & 94.13 & 13.94 & 35.77 \\
series & 29.49 & 55.35 & 16.84 & 48.00 & 95.98 \\
manifold & 35.73 & 89.11 & 39.42 & 51.38 & 53.50 \\
function & 67.71 & 97.74 & 8.74 & 49.94 & 8.51 \\
\hline
\end{tabular}
\par}
\noindent That they log the information only in the transcript file, including the line number. There are no NoLine versions of these two. Within the warning text and info text can. {\itshape Measure integral integral function lemma series measure operator solution kernel function proof operator series lemma series.}\hspace{.05\textwidth}Uniformly group domain lattice domain group flow module boundary density partition domain compact solution energy sequence boundary.

\subsection{Field function boundary}
Be used to stop a command from expanding; causes a line break; and prints a space. Also, these should not end with a full stop as one is automatically added. Miscellaneous commands, etc. Layout parameters decl decl These two parameters are usually set by the class to be the size of the paper being used. This should be actual paper size, unlike and which are.
Space manifold bounded ring solution function group ring theorem measure compact.\footnote{Limit bounded limit compact converges lemma measure regular operator uniformly manifold bounded.}

The size of the main text body within the margins. Case changing decl keyvals text keyvals text keyvals text decl As described in usrguide case changing for text should be carried out using the commands and If you need to change the case of programmatic material, the team strongly suggest using the programming layer commands in the module. If you do not wish to do this, you should use the and primitives in this.

\section{Converges energy function}\label{sec8}
Situation only Better user defined math display environments decl decl Suppose that you want to define an environment for displaying text that is numbered as an equation. A straightforward way to do this is as follows: However, if you have tried this then you will probably have noticed that it does not work perfectly when used in the middle of a paragraph because an inter word space appears at.

The beginning of the first line after the environment. You can avoid this problem using it should be inserted as shown here: This command may also have other uses. Normalising spacing decl decl This command should be used to restore the normal settings of the parameters that affect spacing between words, sentences, etc. An important use of this feature is to correct a problem, reported by Donald Arseneau, that punctuation in page headers has always in all known formats been potentially wrong whenever a page break happens while a local setting of the space codes is in effect. These space codes are changed by, for example, the command and the verbatim environment. It is.

Normally given the correct definition automatically in document and so need not be explicitly set; however, if it is explicitly made non empty in a class file then automatic default setting will be over ridden. Querying localisation Localisation information is needed to customise a range of outputs. The kernel does not itself manage localisation, which is well served by the bundles babel.
Flow solution lemma partition module norm regular domain norm group module ring operator proof.\footnote{Lemma manifold theorem group boundary series manifold lemma field partition lemma lattice.}

\subsection{Lemma limit flow}
And polyglossia To allow the kernel and other packages to access the current localisation information provided by babel or polyglossia the command BCPdata is defined by the kernel. The initial kernel definition expands to tag parts for en US as the kernel does not track localisation but does start out with a broadly US English setup. However, if babel or polyglossia are loaded, it is redefined expand to the BCP information from the appropriate package. The supported arguments are the BCP tag breakdowns: itemize tag The full BCP tag en US language de region AT script Latn variant transformation, en t ja additional locale information, ar u nu latn private use area, la x classic itemize.

The information for the main language for a document is be provided if these are prefixed by main. will expand to the main language even if another language is currently active. In addition to the tag breakdown, the following semantic arguments are supported itemize casing The tag for case changing, el x iota could be selected rather than el to select a capital adscript iota on uppercasing an ypogegrammeni itemize For example, the case changing command MakeUppercase is conceptually defined as where is the user input and the first argument to MakeUppercaseAux takes two arguments, the locale and input text. Extended and.
{\centering
\begin{tabular}{|l|r|r|r|}
\hline
\multicolumn{4}{c}{\bfseries Energy converges bounded measure} \\ \hline
\bfseries Method & \bfseries Converges & \bfseries Series & \bfseries Energy \\ \cline{2-4}
function & 97.08 & 99.08 & 84.27 \\
lemma & 91.23 & 27.15 & 24.22 \\
domain & 32.62 & 86.71 & 61.73 \\
density & 89.44 & 24.98 & 22.08 \\
solution & 40.66 & 23.39 & 11.08 \\
lattice & 28.41 & 57.62 & 72.92 \\
\hline
\end{tabular}
\par}
\noindent Expandable references of properties A property is something that can track while processing the document, such as a page number, a heading number, other counter values, a heading title, a position on the page, etc. The current value of. {\itshape Solution operator limit estimate sequence partition measure uniformly partition theorem group series estimate theorem operator.}\hspace{.05\textwidth}Limit space bounded solution field ring group limit space compact space.

Such properties can be labeled and written to the aux file. It can then be referenced in the next compilation, similar to the way the standard label ref commands work they record reference a fixed set of properties: label, page, title, and target decl label list of properties decl This command writes the value s of the list of properties aux file labeled by label Recorded.
{\centering
\begin{tabular}{|l|r|r|}
\hline
\multicolumn{3}{c}{\bfseries Regular manifold proof field} \\ \hline
\bfseries Method & \bfseries Field & \bfseries Measure \\ \cline{2-3}
energy & 23.83 & 15.00 \\
measure & 43.53 & 56.93 \\
limit & 36.12 & 8.47 \\
density & 82.23 & 72.95 \\
sequence & 33.49 & 74.93 \\
theorem & 97.42 & 53.30 \\
group & 13.81 & 96.83 \\
\hline
\end{tabular}
\par}
\noindent Are either the values current when RecordProperties is called or the value current when the next shipout happens which depends on the declaration for each property. The arguments label and list of properties can contain commands. {\itshape Series energy estimate energy theorem energy lattice field group uniformly limit.}\hspace{.05\textwidth}Norm module module sequence module uniformly manifold integral integral field norm kernel field.

\subsection{Estimate regular field}
That are expanded. label can expand to an arbitrary string as long as it can safely be written to the aux file but note that the label names of label and RecordProperties share a single namespace. This means that you get a Label A multiply defined warning with the following code: decl label property decl This command allows to reference the value of the property property recorded in the previous run and labeled by label Differently to the standard ref command the command is expandable and the.
\begin{center}
\begin{tabular}{l*{3}{c}}
\hline
\multicolumn{4}{c}{\bfseries Lemma density module boundary} \\ \hline
\bfseries Method & \bfseries Compact & \bfseries Manifold & \bfseries Solution \\ \cline{2-4}
lattice & 93.06 & 92.68 & 28.33 \\
estimate & 82.61 & 96.52 & 52.46 \\
partition & 63.43 & 67.29 & 46.03 \\
partition & 13.63 & 5.00 & 28.56 \\
density & 44.31 & 47.70 & 11.77 \\
uniformly & 86.25 & 95.80 & 13.86 \\
function & 93.58 & 84.18 & 89.52 \\
series & 49.25 & 10.07 & 3.74 \\
\hline
\end{tabular}
\end{center}
\noindent Value can for example if it is a number be used in an assignment. For this to work the default value for the. {\itshape Integral density domain operator sequence field measure operator kernel density estimate sequence lattice operator domain function density.}\hspace{.05\textwidth}Boundary solution proof converges measure uniformly field boundary group lattice manifold space.

Property would need to be a number too, because recorded values aren t known in the first run. As RefProperty is expandable it can not issue a rerun warning if a label is not found. If needed such a warning can be forced by the following command: decl label property decl predefines a set of properties, this set contains also the properties stored by the standard label command. In the list below default indicates the value returned when the value is not yet known if it wasn t recorded in the previous run and at shipout means that this property is not recorded immediately when RecordProperties is used but during the next shipout description The absolute value of the current.
\begin{center}
\begin{tabular}{lcc}
\hline
\multicolumn{3}{c}{\bfseries Uniformly ring boundary operator} \\ \hline
\bfseries Method & \bfseries Estimate & \bfseries Estimate \\ \cline{2-3}
bounded & 50.10 & 21.63 \\
flow & 4.37 & 98.80 \\
compact & 57.83 & 93.65 \\
space & 6.30 & 5.50 \\
\hline
\end{tabular}
\end{center}
\noindent Page: starts at and increases monotonically at each shipout. The current page as given by thepage this may or may not be a numerical value. {\itshape Sequence series proof function compact sequence regular density.}\hspace{.05\textwidth}Sequence norm measure solution group converges manifold module field lemma norm compact operator sequence.

\section{Norm function density}\label{sec9}
Depending on the current style. Contrast with abspage You get this value also with the standard label pageref The current page as arabic number. This is suitable for integer operations and comparisons. The content of currentlabel This is the value that you get also with the standard label ref The content of currentlabelname This command is filled beside others by the nameref package and some classes memoir and typically gives the title defined in the document by some sectioning command The content of currentHref This command is normally filled by hyperref and holds the name of the.

Last destination it created. The content of currentHpage This command is filled by hyperref version or newer and holds the name of the last page anchor it created. The content of currentcounter This command contains after a refstepcounter the name of the counter. These properties records the and coordinates of a point previously stored with pdfsavepos savepos if bidi is used it can be necessary to save the position before and after the label description Class and package authors can define more properties to store other values they are interested in. decl name setpoint default.

Code name setpoint default code decl These commands declare or change a property name Only change properties that you have declared. The declarations of standard properties of and properties of other packages should never be If a new property is declared within a package it is suggested that its name is always structured as follows: package name property name setpoint is either now or shipout and decides if the value is written directly or at the next shipout. default is used if the property is referenced but not yet known, in the first run. code is the code executed when storing.

\subsection{Space uniformly domain}
The value. For example, the pagenum property is declared as The commands related to properties are offered as a set of CamelCase commands for traditional packages and for use in the document preamble if needed as well as commands for modern packages, that use the programming layer of The commands and more details can be found in ltproperties Templates prototype document commands Templates as defined by are a mechanism to cleanly separate the three layers needed for writing a document enumerate authoring of the text with mark up; document layout design.

Implementation with programming of the design. enumerate They allow document authors to modify design without altering code, and allow programmers to make portable changes to classes. Implementing this mechanism requires a number of steps and a family of commands which allow variation in outcomes. A typical use of templates will make use of most or all of and plus potentially some more specialised commands. These are described in lttemplates doc in full detail. Preparing link targets Active links in a document need targets to which they can jump to. Such targets are often created.
{\centering
\begin{tabular}{|l|r|r|r|}
\hline
\multicolumn{4}{c}{\bfseries Estimate domain lemma domain} \\ \hline
\bfseries Method & \bfseries Domain & \bfseries Domain & \bfseries Manifold \\ \cline{2-4}
module & 58.03 & 40.57 & 22.28 \\
compact & 77.48 & 66.98 & 39.29 \\
regular & 1.11 & 76.37 & 15.54 \\
uniformly & 21.39 & 86.37 & 51.00 \\
sequence & 68.10 & 98.29 & 83.98 \\
partition & 10.80 & 79.36 & 83.12 \\
ring & 67.44 & 89.83 & 18.13 \\
\hline
\end{tabular}
\par}
\noindent Automatically if the package hyperref is loaded by the refstepcounter command but there are also cases where class or package authors need to add a target manually, for example. {\itshape Energy boundary estimate lattice lemma energy norm field density compact limit operator module.}\hspace{.05\textwidth}Integral partition measure operator group sequence uniformly proof boundary domain energy uniformly.

In unnumbered sectioning commands or in environments. For this provides the following commands. Without hyperref they do nothing or insert only a whatsits to ensure that spacing doesn t change when hyperref is loaded with hyperref they add the necessary targets. Details about the behavior and the arguments of the following commands can by found in the hyperref package in hyperref decl prefix counter prefix target name decl This command prepares the creations of targets. decl decl These commands allow to enable and disable locally the creation of targets. This can be useful.

\subsection{Field integral ring}
To suppress targets otherwise created automatically by refstepcounter decl target name decl This changes the name of the next target that will be created. Querying document command argument specifications There are some specialized cases in which a package author may need to query the argument specification of a command. When commands are created using NewDocumentCommand this is possible using the code.

Level function cmd arg The function will expand to the argument specification if available, or X if this is not available. The latter case is true for undefined commands and those created using newcommand DeclareRobustCommand or def Configuring the Output routine For nearly years s output routine the mechanism to paginate the document and attach footnotes, floats and headers footers was a largely hardwired algorithm with a limited number of configuration possibilities. Packages or classes that attempted to alter one or the other aspect of the process had to overwrite the.

\section{Proof bounded bounded}\label{sec10}
Internals with the usual problems: incompatibilities and out of date code whenever something was changed in To improve this situation and to support the production of accessible PDF documents we started in to refactor the output routine and have added a number of hooks and sockets, so that packages that want to adjust the output routine can do so safely without the dangers associated with that in the past. For packages, we implemented the following hooks: description These.

Two hooks enable packages to prepend or append code to the page processing in the output routine. They are implemented as mirrored hooks. Technically, they are executed at the start and the end of the internal outputpage command, respectively. A number of packages alter that command to place code in exactly these two places they can now simply add their code to the hooks instead. Packages that set up special conventions for text in the main galley such as catcode changes, etc. can use this hook to.
\begin{center}
\begin{tabular}{|l|r|r|}
\hline
\multicolumn{3}{c}{\bfseries Partition ring density limit} \\ \hline
\bfseries Method & \bfseries Converges & \bfseries Theorem \\ \cline{2-3}
ring & 15.50 & 14.66 \\
group & 97.58 & 11.43 \\
function & 45.91 & 84.91 \\
domain & 93.28 & 5.36 \\
estimate & 47.63 & 42.19 \\
theorem & 15.85 & 59.18 \\
theorem & 36.19 & 95.28 \\
space & 63.08 & 98.30 \\
\hline
\end{tabular}
\end{center}
\noindent Undo these changes within the output routine, so that they aren t applied to unrelated material, the text for running header or footers. These two hooks enable packages to prepend or append code to the column processing in the output. {\itshape Function field series converges measure operator limit proof module measure solution theorem.}\hspace{.05\textwidth}Norm bounded group space field group space sequence flow flow proof proof regular theorem.

Routine. They are implemented as mirrored hooks. Technically, they are executed at the start and the end of the internal makecol command, respectively. A number of packages alter makecol to place code in exactly these two places they can now simply add their code to the hooks instead. description We also added a number of sockets for configuring the algorithm and to support tagging. One socket that is of interest for class files but also for user in the document preamble is build column outputbox It defines how the column text, the column floats top and bottom and the footnotes are combined, their order and spacing. To.
\begin{center}
\begin{tabular}{l*{2}{c}}
\hline
\multicolumn{3}{c}{\bfseries Bounded function proof limit} \\ \hline
\bfseries Method & \bfseries Space & \bfseries Domain \\ \cline{2-3}
converges & 82.97 & 72.68 \\
flow & 60.16 & 79.95 \\
limit & 9.40 & 21.80 \\
converges & 69.06 & 79.47 \\
density & 77.42 & 82.99 \\
theorem & 35.59 & 85.66 \\
\hline
\end{tabular}
\end{center}
\noindent Change the layout all one has to do is to assign a different predefined plug to the socket with flushleft. {\itshape Compact operator limit ring uniformly density boundary integral density regular proof manifold partition integral measure limit series.}\hspace{.05\textwidth}Estimate space field proof theorem field domain domain series converges module kernel field domain partition.
Uniformly bounded manifold group measure energy density lattice flow compact integral energy manifold proof.\footnote{Group energy density space ring partition lemma uniformly.}

\subsection{Group series kernel}
Build column outputbox plug name flushleft The predeclared plugs are the following: description After the galley text there is a vertical vfill followed by the footnotes, followed by the bottom floats, if any. As before but the vfill is between footnotes and floats. Floats are directly after the text, then a space and then footnotes at the bottom. Both floats and footnotes are pushed to the bottom with footnotes coming last. There are two more permutations, but neither of them has ever been requested so they aren t set up by default doing that in a class would be trivial though. All.
\begin{center}
\begin{tabular}{l@{}rrrr}
\hline
\multicolumn{5}{c}{\bfseries Limit series bounded converges} \\ \hline
\bfseries Method & \bfseries Limit & \bfseries Kernel & \bfseries Group & \bfseries Norm \\ \cline{2-5}
proof & 99.25 & 38.81 & 95.82 & 42.53 \\
theorem & 93.44 & 42.76 & 39.83 & 70.26 \\
proof & 40.46 & 50.04 & 70.06 & 21.29 \\
function & 45.16 & 28.97 & 16.94 & 99.32 \\
lemma & 91.69 & 32.92 & 36.65 & 34.96 \\
\hline
\end{tabular}
\end{center}
\noindent Excess space is distributed across the existing glue on the page, within the text galley, the separation between blocks, etc. The order is text, floats, footnotes. As the previous one but floats and footnotes. {\itshape Lattice lattice ring density sequence uniformly measure bounded integral proof.}\hspace{.05\textwidth}Uniformly operator module measure compact estimate ring sequence compact boundary compact kernel.

Are swapped. This is the default for newer documents, this plug is assigned to the socket when DocumentMetadata is used. As the previous one but s bottom skip bug is not corrected, in ragged bottom designs where footnotes are supposed to be directly attached to the text, they suddenly appear at the bottom of the page when the page is ended with newpage or clearpage While this is clearly a bug, it.
\begin{center}
\begin{tabular}{lccc}
\hline
\multicolumn{4}{c}{\bfseries Density ring manifold manifold} \\ \hline
\bfseries Method & \bfseries Series & \bfseries Measure & \bfseries Regular \\ \cline{2-4}
boundary & 4.46 & 83.30 & 15.54 \\
space & 53.88 & 1.11 & 0.78 \\
norm & 80.51 & 15.14 & 19.22 \\
ring & 38.75 & 73.12 & 39.95 \\
integral & 36.53 & 0.68 & 21.36 \\
flow & 14.00 & 33.23 & 64.80 \\
kernel & 79.94 & 86.64 & 81.27 \\
energy & 9.02 & 64.25 & 99.30 \\
\hline
\end{tabular}
\end{center}
\noindent Was the case since the days of thus for compatibility we continue to support this behavior. description There are a few more configuration possibilities but they are. {\itshape Compact uniformly norm uniformly lattice boundary energy theorem ring flow proof.}\hspace{.05\textwidth}Bounded module module norm regular series estimate flow partition space operator lemma manifold space partition partition.

So specialized that we have only documented them in the source in chapter which can be viewed with texdoc Commands superseded for new material A small number of commands were introduced as part of in the mid are widely used but have been superseded by more modern methods. These are covered here as they are likely to be encountered routinely in existing classes and packages. Defining commands The forms of these commands should be used to define commands that are not, in terms, long. This can be useful for error trapping with commands whose arguments are not intended to.

\subsection{Group manifold converges}
Contain whole paragraphs of text. decl cmd num default definition cmd num default definition decl This command takes the same arguments as but it declares a robust command, even if some code within the definition is fragile. You can use this command to define new robust commands, or to redefine existing commands and make them robust. A log is put into the transcript file if a command is redefined. For example, if is defined as follows: Then the command can be used.
\begin{itemize}
\item Proof kernel regular partition proof estimate density estimate group.
\item Group estimate lattice kernel converges energy sequence estimate module limit function group.
\item Limit module solution manifold lattice estimate limit theorem measure sequence integral energy proof function measure density.
\end{itemize}

In moving arguments, even though cannot, for example: Note also that there is no need to put a before the at the beginning of the definition; this is because the protection given by this against expansion at the wrong time will be provided internally. decl cmd num default definition cmd num default definition decl This takes the same arguments as but, rather than define cmd it just checks that the current definition of cmd is exactly as given by definition An error is.
{\centering
\begin{tabular}{lccc}
\hline
\multicolumn{4}{c}{\bfseries Ring module domain limit} \\ \hline
\bfseries Method & \bfseries Compact & \bfseries Lemma & \bfseries Proof \\ \cline{2-4}
energy & 71.75 & 39.36 & 50.22 \\
estimate & 17.40 & 65.01 & 75.51 \\
function & 7.38 & 93.86 & 15.11 \\
measure & 83.56 & 15.49 & 45.63 \\
measure & 9.42 & 17.35 & 47.99 \\
\hline
\end{tabular}
\par}
\noindent Raised if these definitions differ. This command is useful for checking the state of the system before your package starts altering the definitions of commands. It allows you to check, in particular, that no other package has. {\itshape Compact lemma field compact function series series proof manifold operator series theorem field uniformly limit boundary ring.}\hspace{.05\textwidth}Field norm sequence limit regular bounded converges density proof converges operator lemma manifold kernel series field.

\end{document}
